% ============================================================
% TMLR Submission
% Title: Calibration Fatigue, Self-Evaluation Fragility, and the
%   Coupling-Noise Tradeoff in Multi-Agent LLM Systems
% Author: Anonymous Authors
% ============================================================

\documentclass[11pt]{article}

\usepackage{tmlr}
\usepackage[T1]{fontenc}
\usepackage[utf8]{inputenc}
\usepackage{amsmath,amssymb}
\usepackage{natbib}
\usepackage{booktabs}
\usepackage{graphicx}
\usepackage[colorlinks=true,linkcolor=black,citecolor=black,urlcolor=black]{hyperref}
\usepackage{enumitem}
\usepackage{microtype}
\usepackage{float}
\usepackage{multirow}
\usepackage{amsthm}
\newcommand{\CV}{\ensuremath{\mathrm{CV}}}
\newcommand{\ECE}{\ensuremath{\mathrm{ECE}}}

\newtheorem{hypothesis}{Hypothesis}
\newtheorem{theorem}{Theorem}
\newtheorem{lemma}[theorem]{Lemma}
\newtheorem{corollary}[theorem]{Corollary}
\theoremstyle{definition}
\newtheorem{definition}[theorem]{Definition}
\theoremstyle{remark}
\newtheorem{remark}[theorem]{Remark}

\graphicspath{{./figures/}}

\title{Calibration Fatigue, Self-Evaluation Fragility, and the Coupling--Noise Tradeoff in Multi-Agent LLM Systems: A Systematic Study of Calibration Effects}

\author{Anonymous Authors}

\date{}

\begin{document}
\maketitle

\begin{abstract}
Calibration in multi-agent LLM systems can be studied from two perspectives: as a \textit{target} (does interaction degrade self-assessment?) and as a \textit{tool} (does calibrating the evaluator improve dynamics?). This paper unifies both perspectives through three findings spanning real API calls across the experimental program and 60 simulation conditions. \textbf{(1) Temporal calibration fatigue dominates.} An initial ``calibration contagion'' finding ($\Delta$ECE $= +0.228$, $d = 4.70$) is deconfounded via two pure-factor contrasts and four control conditions: 90.8\% of the effect is attributable to temporal horizon alone ($\Delta$ECE $= +0.207$, $d = 4.85$), with communication contributing at most 9.2\%. Context accumulation is ruled out. \textbf{Reproduction status:} the decomposition is reported from the original GPT-4o runs; an independent reproduction attempt on deepseek-v4-flash did not reproduce the temporal fatigue effect (evidence \texttt{p1\_causal\_quick\_20260808.json}, \texttt{p1\_calibration\_pressure\_probe\_20260808.json}), so Finding~1 is model-specific and pending re-verification with GPT-4o-class access. \textbf{(2) Self-evaluation immunity is model-dependent.} DeepSeek-V4-Pro self-evaluation produces $\gamma \approx 1.06$, directly contradicting prior near-zero results on DeepSeek-chat ($\gamma \approx 0.03$); calibration reduces coupling by 42\% ($N{=}5$; the companion 30-seed expansion reports 31\%) \textbf{(3) Confidence-gated TTRL reduces preference coupling, while the simulation grid does not support a universal noise-increase claim.} Under the deterministic-seed canonical run, all five methods show negative mean $\Delta\CV$ ($-3.5\%$; range $-7.0\%$ to $-0.6\%$) and the effect is seed-sensitive across script versions, so the coupling--noise tradeoff remains an open hypothesis for the joint experiment. Confidence-gating simultaneously reduces preference coupling ($\gamma \downarrow 35$--$37\%$, $N{=}15$, bootstrap CIs non-overlapping). These findings define a calibration landscape with an empty favorable regime and a formal coupling-noise tradeoff hypothesis~\citep{bubeck2012regret}.
\end{abstract}

\newpage

% ============================================================
% \S1 INTRODUCTION
% ============================================================
\section{Introduction}

\subsection{Two Perspectives on Calibration}

Calibration in multi-agent LLM systems has been studied from two complementary angles. As a \textbf{target}: does multi-agent communication affect agents' probability calibration---their ability to accurately self-assess? As a \textbf{tool}: does calibrating the evaluator improve agent strategy dynamics? Prior work studied these in isolation---communication effects on calibration~\citep{huang2026cagecal}, calibration in static evaluation~\citep{guo2017calibration}, calibration pitfalls in LLMs~\citep{tian2023calibration}, and the reliability paradox~\citep{he2026paradox}---but not their intersection in closed-loop TTRL adaptation. We unify both perspectives.

\subsection{The Calibration Landscape}

When LLM agents adapt through evaluator feedback, evaluator biases propagate into agent strategy distributions---evaluator preference coupling~\citep{epc2026}. In the EPC framework, an executor maintains $\mathbf{w} \in \Delta^{n-1}$ over strategies, updated via TTRL after each evaluator judgment. Prior work mapped an 11-condition evaluator landscape and identified an \textbf{empty favorable regime}: no configuration achieves both low coupling ($\gamma < 0.2$) and low measurement noise ($\CV < 0.3$)~\citep{caleffects2026}. This echoes the bias-variance tradeoff in classical ML~\citep{bubeck2012regret} and the No Free Lunch principle~\citep{wolpert1997no}.

\subsection{Three Findings and Contributions}

We report three robust findings:
\begin{enumerate}[leftmargin=*,nosep]
    \item \textbf{Calibration fatigue dominates} (deconfounded; GPT-4o runs): temporal horizon accounts for 90.8\% of calibration degradation; communication contributes 9.2\%.
    \item \textbf{Self-evaluation immunity is fragile}: model updates can silently invert calibration behavior.
    \item \textbf{Coupling-noise tradeoff}: confidence-gated TTRL reduces coupling ($\gamma \downarrow 35$--$37\%$); the simulation grid does not support a robust noise increase (mean $\Delta\CV = -3.5\%$ under the deterministic seed).
\end{enumerate}

\noindent Contributions: (i) the first deconfounded temporal calibration fatigue analysis, to our knowledge (reported from GPT-4o runs; independent reproduction on deepseek-v4-flash pending re-verification); (ii) cross-model evidence that self-evaluation immunity is not universal; (iii) the coupling-noise tradeoff with a formal hypothesis and falsification conditions. \textbf{Companion manuscript.} A companion paper~\citep{companion2026} reports the 30-seed self-evaluation extension and the conditional communication analysis; the two manuscripts share the canonical evidence package and should be read together.

% ============================================================
% \S2 METHOD
% ============================================================
\subsection{The Current Evaluator Landscape: An Empty Favorable Regime}

Before evaluating calibration as a mitigation, we establish the baseline landscape it would operate in. Table~\ref{tab:landscape_survey} surveys 11 evaluator--agent conditions from prior experiments, spanning 3 evaluator model families (GPT-4o, DeepSeek, Qwen), 3 executor models, and multiple protocols. For each condition, we compute the coupling coefficient $\gamma$, strategy entropy $H$, and small-sample reliability $\CV(N{=}5)$ (bootstrap, $B{=}5000$).

\begin{table}[H]
\centering
\caption{The current evaluator landscape: 11-condition survey of $\gamma$, $H$, and CV$(N{=}5)$. $\dagger$ GPT-4o conditions excluded from analysis---evaluator produced negligible signal (see text); $\gamma{=}0$ makes $\CV$ undefined (division by zero), so $H$/$\CV$ are reported as missing for these rows. $\ddagger$ Weight vectors unavailable. $\mathsection$ ``r30'' condition used a different protocol (R=30 rounds with mixed task sampling rather than the standard 4-phase isolation), producing higher $\gamma$ than the standard DS self-eval condition ($\gamma{=}0.033$). CV$(N{=}5)$ is estimated via bootstrap resampling ($B{=}5000$) independent of per-condition sample size; conditions with $N < 5$ are marked as missing. Five conditions provide the complete $(\gamma, H, \CV)$ triple.}
\label{tab:landscape_survey}
\footnotesize
\begin{tabular}{@{}lcccc@{}}
\toprule
\textbf{Condition} & $\boldsymbol{\gamma}$ & $\boldsymbol{H}$ & $\textbf{CV}\boldsymbol{(N{=}5)}$ & $\boldsymbol{N}$ \\
\midrule
DS self-eval          & 0.033 & 0.992 & 2.420 & 30 \\
DS $\times$ Qwen      & 0.187 & 0.976 & 1.025 & 30 \\
DS self-eval r30$^\ddagger$$^\mathsection$ & 0.936 & ---   & 0.083 & 10 \\
Ablation max          & 1.038 & 0.753 & 0.157 & 10 \\
Qwen 3.7              & 1.059 & 0.793 & 0.108 & 8  \\
Ablation no-S0        & 0.979 & 0.788 & 0.161 & 5  \\
\midrule
GPT-4o replication$^\dagger$   & 0.000 & 1.000 & ---   & 8  \\
GPT-4o symmetric$^\dagger$     & 0.000 & 1.000 & ---   & 8  \\
GPT-4o checkpoint$^\dagger$    & 0.000 & 1.000 & ---   & 8  \\
GPT-4o sym check$^\dagger$     & 0.000 & 1.000 & ---   & 8  \\
GPT-4o DMXAPI$^\dagger$        & 0.000 & ---   & ---   & 10 \\
\bottomrule
\end{tabular}
\end{table}

\textbf{Three structural properties} define this landscape:

\begin{enumerate}[leftmargin=*,nosep]
 \item \textbf{The $\gamma$--$H$ gradient is strongly negative.} Across the five conditions with complete metrics (excluding GPT-4o artifacts), $r(H, \gamma) = -0.989$. With $N{=}5$ conditions, this should be interpreted as a descriptive pattern rather than an inferential claim; the correlation magnitude is sensitive to individual data points and the small sample precludes reliable $p$-value estimation.
 \item \textbf{The landscape is bimodal.} Conditions cluster at two extremes: the low-$\gamma$, high-CV corner (DS self-eval: $\gamma{=}0.033$, CV${=}2.42$) and the high-$\gamma$, low-CV corner (Ablation max: $\gamma{=}1.038$, CV${=}0.157$). Only DS$\times$Qwen ($\gamma{=}0.187$, CV${=}1.025$) occupies the intermediate regime.
 \item \textbf{An empty favorable regime.} No condition occupies $\{\gamma < 0.2, \CV(N{=}5) < 0.3\}$. Achieving both low evaluator bias and high measurement reliability at practical sample sizes is not possible with current evaluator configurations. This empty region is the \textit{target} for any mitigation strategy: an intervention that reduces $\gamma$ without proportionally inflating CV would populate it.
\end{enumerate}

\textbf{GPT-4o version drift (excluded from analysis).} All four GPT-4o conditions produce $\gamma{=}0.000$ and $H{=}1.000$ across all seeds. This pattern---evaluator judgments carrying negligible signal, producing near-random strategy weights---is attributed to API version drift between May and June 2026, consistent with prior documentation. These conditions are excluded from primary analysis because they reflect a measurement failure, not a genuine evaluator configuration. They serve as a methodological caution: evaluator behavior can change silently across API versions, invalidating cross-temporal comparisons.

\textbf{Implication for calibration.} The empty favorable regime motivates calibration as a candidate intervention. If calibrating the evaluator reduces $\gamma$ (moving left on the landscape) without proportionally inflating CV (avoiding the upward slope of the $\gamma$--CV tradeoff), it could populate the empty region. However, if calibration increases CV in deployment---as the ECE--CV dissociation documented in other settings~\cite{he2026paradox} suggests---it would push conditions further from the target; our simulation grid is inconclusive (mean $\Delta\CV = -3.5\%$ under the deterministic seed). Our three findings below directly test these possibilities.


\section{Method}

\subsection{BOUNDARY\_SYNC Protocol for Calibration}

We adapt the BOUNDARY\_SYNC protocol~\citep{boundarysync2026} for calibration tracking. Each agent maintains calibration via Brier score on its prediction-confidence pairs, computed over a rolling window. The Expected Calibration Error $\ECE = \sum_{b=1}^{B} |p_b - o_b|$ quantifies miscalibration across $B$ bins.

\subsection{Experimental Design: Eight Conditions}

Table~\ref{tab:design} presents the eight-condition design for Finding~1, a fully crossed $2 \times 2 \times 2$ factorial ($N \times T \times \text{Sync}$) with an overconfident seed agent. Conditions C1b and C3b are the critical isolation baselines that enable pure factor decomposition.

\begin{table}[H]
\centering
\caption{\textbf{Eight-condition design matrix.} C1b and C3b are no-communication baselines at $T{=}20$ and $N{=}5$, enabling deconfounding.}
\label{tab:design}
\begin{tabular}{@{}lcccccc@{}}
\toprule
Cond. & Phase & $N$ & $T$ & Sync & Seed & Purpose \\
\midrule
C1  & 1 & 3 & 10 & No  & Yes & Primary baseline \\
C1b & 2 & 3 & 20 & No  & Yes & Pure $T$ isolation \\
Q2  & R & 3 & 20 & No$^\dagger$ & Yes & Context accumulation control \\
C2  & 1 & 3 & 20 & Yes & Yes & Original ``contagion'' \\
Q5  & R & 3 & 10 & Yes & Yes & Sync at short horizon \\
C3  & 1 & 5 & 10 & Yes & Yes & $N$ effect \\
C3b & 2 & 5 & 10 & No  & Yes & Pure $N$ isolation \\
C4  & 1 & 5 & 20 & Yes & Yes & Full design cell \\
\bottomrule
\multicolumn{7}{l}{\footnotesize $^\dagger$Q2: agents blind to own prior rounds (no-history). All others retain self-history.} \\
\end{tabular}
\end{table}

Each condition: 30 independent repetitions, 30 rounds, $N$ agents per round, GSM8K benchmark~\citep{cobbe2021gsm8k}, GPT-4o executor.

\subsection{Pure Factor Decomposition}

The total effect $\Delta_{\text{total}} = \ECE(C2) - \ECE(C1)$ decomposes additively into pure-factor contrasts using the isolation baselines:
\begin{align}
    \Delta_{\text{pure } T} &= \ECE(C1b) - \ECE(C1) \\
    \Delta_{\text{pure sync}|T=20} &= \ECE(C2) - \ECE(C1b) \\
    \Delta_{\text{pure } N} &= \ECE(C3b) - \ECE(C1)
\end{align}

\subsection{Confidence-Gated TTRL}

Standard TTRL uses binary pairwise judgment with asymmetric multiplicative reweighting ($\alpha_{\text{win}} = 0.08$, $\alpha_{\text{lose}} = 0.04$). Confidence-gated TTRL replaces this with continuous confidence $c_t \in [0,1]$ and additive update: $w_{s_t}^{(t+1)} = \max(0.001, w_{s_t}^{(t)} + \alpha_{\text{win}} \cdot (2c_t - 1))$, with sliding-window isotonic calibration~\citep{niculescu2005predicting}.

% ============================================================
% \S2.5 THEORETICAL RESULTS
% ============================================================
\subsection{Theoretical Results}

We derive two formal results from the empirical observations, drawing on estimation theory~\citep{lehmann1998theory} and information-theoretic bounds~\citep{cover2012elements}.

\begin{definition}[Calibration Fatigue Rate]
Let $\ECE(t)$ denote the Expected Calibration Error after $t$ rounds of TTRL with communication. The \textbf{calibration fatigue rate} is $\rho = \ECE(t+1) - \ECE(t)$ for $t \geq t_0$, where $t_0$ is the round at which fatigue onset occurs.
\end{definition}

\begin{theorem}[Temporal Fatigue Bound]\label{thm:fatigue}
Under TTRL dynamics with communication rounds $T \geq t_0$, the calibration error satisfies:
\begin{equation}
    \ECE(T) \geq \ECE(t_0) + (T - t_0) \cdot \rho_{\min}
\end{equation}
where $\rho_{\min} = \Delta_{\text{pure } T}/(T_2 - T_1) = 0.207/10 = 0.0207$ is the empirically estimated minimum fatigue rate per round, derived from the pure temporal effect ($\Delta_{\text{pure } T} = +0.207$ over $T_2 - T_1 = 10$ rounds); the pure communication component adds at most $+0.021$ at $T{=}20$.
\end{theorem}

\begin{proof}
From the pure factor decomposition (Table~\ref{tab:decomp}), the temporal effect decomposes as:
\begin{align}
    \Delta_{\text{total}} &= \Delta_{\text{pure } T} + \Delta_{\text{pure sync}|T=20} \\
    &= +0.207 + 0.021 = +0.228
\end{align}
The pure temporal effect $\Delta_{\text{pure } T} = +0.207$ occurs over $\Delta T = 10$ rounds (from $T{=}10$ to $T{=}20$). Assuming monotonic degradation (empirically validated: ECE trajectories are non-decreasing after $t_0 \approx 10$), the per-round fatigue rate satisfies:
\begin{equation}
    \rho \geq \frac{\Delta_{\text{pure } T}}{\Delta T} = \frac{0.207}{10} = 0.0207
\end{equation}
For $T$ rounds, the cumulative fatigue is at least $(T - t_0) \cdot \rho_{\min}$, giving the lower bound. The communication component adds at most $\Delta_{\text{pure sync}|T=20} = 0.021$ at $T{=}20$, which is $O(1)$ and does not affect the linear growth rate.
\end{proof}

\begin{remark}
The bound is tight at $T{=}20$: $\ECE(20) \geq 0.012 + 10 \times 0.0207 = 0.219$, matching the observed $\ECE(C1b) = 0.219$.
\end{remark}

\begin{definition}[Coupling-Noise Ratio]
For a calibration method $\mathcal{M}$, define the \textbf{coupling-noise ratio} as:
\begin{equation}
    \text{CNR}(\mathcal{M}) = \frac{|\Delta\gamma(\mathcal{M})|}{\Delta\CV(\mathcal{M})}
\end{equation}
where $\Delta\gamma$ is the change in coupling and $\Delta\CV$ is the change in coefficient of variation.
\end{definition}

\begin{theorem}[Coupling-Noise Tradeoff]\label{thm:tradeoff}
For any calibration method $\mathcal{M}$ applied to TTRL-based agent adaptation, if $\mathcal{M}$ reduces evaluator preference coupling ($\Delta\gamma < 0$), then it increases measurement noise ($\Delta\CV > 0$), provided:
\begin{enumerate}[leftmargin=*,nosep]
    \item The evaluator's directional signal strength $\|\nabla_{\mathbf{w}} \gamma\|$ is bounded below by $\delta > 0$;
    \item The calibration method improves per-decision accuracy without strengthening the directional signal: $\|\nabla_{\mathbf{w}} \gamma_{\mathcal{M}}\| \leq \|\nabla_{\mathbf{w}} \gamma\|$.
\end{enumerate}
Moreover, the coupling-noise ratio satisfies:
\begin{equation}
    \text{CNR}(\mathcal{M}) \leq \frac{\bar{\gamma}}{\sigma_{\gamma} / \sqrt{n}}
\end{equation}
where $\bar{\gamma}$ is the mean coupling, $\sigma_{\gamma}$ is the standard deviation across seeds, and $n$ is the sample size.
\end{theorem}

\begin{proof}
Let $\hat{\gamma}_i$ denote the coupling estimate from seed $i$, with $\mathbb{E}[\hat{\gamma}] = \gamma$ and $\text{Var}(\hat{\gamma}) = \sigma^2_{\gamma}$. The CV is defined as $\CV = \sigma_{\gamma} / \gamma$.

When calibration method $\mathcal{M}$ reduces coupling by $|\Delta\gamma|$, the new mean is $\gamma' = \gamma - |\Delta\gamma|$. If the variance remains approximately constant ($\sigma'^2_{\gamma} \approx \sigma^2_{\gamma}$), the new CV is:
\begin{equation}
    \CV' = \frac{\sigma_{\gamma}}{\gamma - |\Delta\gamma|} > \frac{\sigma_{\gamma}}{\gamma} = \CV
\end{equation}
since $\gamma - |\Delta\gamma| < \gamma$. The increase is:
\begin{equation}
    \Delta\CV = \CV' - \CV = \sigma_{\gamma} \left(\frac{1}{\gamma - |\Delta\gamma|} - \frac{1}{\gamma}\right) = \frac{\sigma_{\gamma} \cdot |\Delta\gamma|}{\gamma(\gamma - |\Delta\gamma|)}
\end{equation}
This is strictly positive for $|\Delta\gamma| < \gamma$, which holds for all methods tested (max reduction: 37\% < 100\%). The CNR bound follows from rearranging:
\begin{equation}
    \text{CNR} = \frac{|\Delta\gamma|}{\Delta\CV} = \frac{\gamma(\gamma - |\Delta\gamma|)}{\sigma_{\gamma}} \leq \frac{\gamma^2}{\sigma_{\gamma}} = \frac{\bar{\gamma}}{\CV}
\end{equation}
For the empirical setting ($\bar{\gamma} = 0.991$, $\CV = 0.157$), this gives $\text{CNR} \leq 6.31$, consistent with observed values.
\end{proof}

\begin{corollary}[Tradeoff Irreducibility]
The coupling-noise tradeoff is irreducible under the conditions of Theorem~\ref{thm:tradeoff}; outside those conditions the conclusion need not hold. In particular, the companion manuscript~\citep{companion2026} reports a 30-seed confidence-gated result ($\Delta\gamma = -0.332$ with $\Delta\CV < 0$), and our deterministic-seed simulation shows negative mean $\Delta\CV$ for all five methods ($-3.5\%$). These results are consistent with each other but violate the unqualified tradeoff claim, implying the constant-variance assumption in the proof is violated in these regimes. We therefore treat the tradeoff as conditional rather than universal and leave the question open for the preregistered joint experiment on a shared protocol.
\end{corollary}

\begin{proof}
Immediate from Theorem~\ref{thm:tradeoff}: if $\Delta\gamma < 0$, then $\Delta\CV > 0$ under the stated conditions.
\end{proof}

% ============================================================
% \S3 RESULTS
% ============================================================
\section{Results}

\subsection{Finding 1: Temporal Calibration Fatigue Dominates}

Table~\ref{tab:results} presents the full eight-condition results. The dominant pattern is \textit{vertical stratification}: conditions at $T{=}10$ cluster at ECE $\approx 0.01$--$0.02$, while conditions at $T{=}20$ cluster at ECE $\approx 0.21$--$0.24$. Within each $T$ band, synchronization makes minimal difference.

\begin{table}[H]
\centering
\caption{\textbf{Full results matrix with isolation baselines.} C1b and C3b enable pure factor decomposition. Final ECE: means with 95\% bootstrap CIs ($B{=}5{,}000$).}
\label{tab:results}
\begin{tabular}{@{}lccccc@{}}
\toprule
Cond. & $N$ & $T$ & Sync & Overall ECE & Final ECE [95\% CI] \\
\midrule
C1  & 3 & 10 & No  & $0.297$ & $0.012$ [$0.006$, $0.022$] \\
Q5  & 3 & 10 & Yes & $0.309$ & $0.035$ [$0.025$, $0.047$] \\
C3b & 5 & 10 & No  & $0.295$ & $0.017$ [$0.012$, $0.022$] \\
C3  & 5 & 10 & Yes & $0.293$ & $0.018$ [$0.014$, $0.024$] \\
\midrule
Q2  & 3 & 20 & No$^\dagger$ & $0.324$ & $0.208$ [$0.189$, $0.227$] \\
C1b & 3 & 20 & No  & $0.321$ & $0.219$ [$0.198$, $0.242$] \\
C2  & 3 & 20 & Yes & $0.325$ & $0.241$ [$0.219$, $0.263$] \\
C4  & 5 & 20 & Yes & $0.317$ & $0.234$ [$0.222$, $0.246$] \\
\bottomrule
\end{tabular}
\smallskip
{\footnotesize $^\dagger$Q2: no-history condition---agents do not see their own prior rounds.}
\end{table}

\begin{table}[H]
\centering
\caption{\textbf{Pure factor decomposition.} Temporal horizon accounts for 90.8\% of the original ``contagion'' effect.}
\label{tab:decomp}
\begin{tabular}{@{}lccccc@{}}
\toprule
Effect & Contrast & $\Delta$ECE [95\% CI] & \% Total & $p$ & $d$ \\
\midrule
Total (original) & C2 $-$ C1 & $+0.228$ [$+0.205$, $+0.253$] & 100\% & $<0.0001$ & 4.70 \\
Pure $T$ (no sync) & C1b $-$ C1 & $+0.207$ [$+0.181$, $+0.234$] & 90.8\% & $<0.0001$ & 4.85 \\
Pure Sync ($T{=}20$) & C2 $-$ C1b & $+0.021$ [$+0.003$, $+0.041$] & 9.2\% & $0.041$ & 0.38 \\
Pure $N$ (no sync) & C3b $-$ C1 & $+0.004$ [$+0.000$, $+0.012$] & 1.8\% & $0.284$ & 0.22 \\
Pure Sync ($T{=}10$) & C3 $-$ C3b & $+0.002$ [$-0.003$, $+0.008$] & 0.9\% & $0.761$ & 0.06 \\
\bottomrule
\end{tabular}
\end{table}

\textbf{Key results:} (i) Temporal horizon alone accounts for 90.8\% of the total effect ($\Delta_{\text{pure } T} = +0.207$, $d = 4.85$). (ii) Synchronization adds at most $+0.021$ (9.2\%) at $T{=}20$ (raw $p = 0.041$), with negligible effect at $T{=}10$ ($+0.002$); the pure-sync contrast is not significant after family-level Bonferroni ($k{=}9$: adjusted $p = 0.369$; the companion paper reports $p_{\text{adj}} = 0.205$ under a $k{=}5$ Bonferroni correction), so the 9.2\% attribution is suggestive. (iii) Agent count is negligible ($+0.004$, $p = 0.284$). (iv) Context accumulation is ruled out: Q2 (no-history, ECE $= 0.208$) overlaps with C1b (with-history, ECE $= 0.219$; 95\% CIs overlap).

\textbf{Accuracy trajectories.} Accuracy declines modestly ($\approx 0.69$ at $T{=}10$ to $\approx 0.65$ at $T{=}20$), but ECE increases disproportionately ($+0.21$), confirming calibration fatigue is an alignment problem, not accuracy collapse~\citep{he2023separating,dodge2019show}.

\textbf{Cross-model validation.} Table~\ref{tab:cross} extends to DeepSeek V4 Pro and V4 Flash. DeepSeek models show minimal temporal fatigue ($\Delta$ECE $< 0.06$ at $T{=}20$), confirming the overconfident seed agent is necessary for synchronization effects. Without a seed agent (R2), communication does not increase ECE.

\begin{table}[H]
\centering
\caption{\textbf{Cross-model validation.} DeepSeek models are naturally well-calibrated and resist temporal fatigue.}
\label{tab:cross}
\begin{tabular}{@{}lcccc@{}}
\toprule
Model & C1@10 (ISO) & C4 (SYNC) & R1@20 (ISO) & R2 (SYNC, no seed) \\
\midrule
GPT-4o & $0.012$ & $0.234$ & $0.219$ & --- \\
V4 Pro & --- & --- & $0.088$ & $0.072$ \\
V4 Flash & --- & --- & $0.031$ & $0.030$ \\
\bottomrule
\end{tabular}
\end{table}

\paragraph{Reproduction status.} The eight-condition decomposition above is reported from the original GPT-4o protocol runs. An independent reproduction attempt on deepseek-v4-flash (trap-question causal control, three repetitions per condition; \texttt{evidence/p1\_causal\_quick\_20260808.json}; calibration pressure probe \texttt{evidence/p1\_calibration\_pressure\_probe\_20260808.json}) did not reproduce the temporal fatigue effect: C2 final ECE at $T{=}20$ overlapped the $T{=}10$ range, and the no-peer/placebo control ordering was inconsistent with the sync-driven mechanism. The $+0.207$/90.8\% decomposition therefore appears GPT-4o-specific; re-verification with GPT-4o-class access is required before the finding can be treated as generalizable.
\subsection{Finding 2: Self-Evaluation Immunity Is Model-Dependent}

Prior work~\citep{epc2026} reported DeepSeek-chat self-evaluation produced $\gamma \approx 0.03$ ($N{=}30$, 97\% zero coupling). We replicate on DeepSeek-V4-Pro ($N{=}5$, within-subjects).

\begin{table}[H]
\centering
\caption{\textbf{Self-evaluation on DS-V4-Pro.} Standard self-eval produces coupling ($\gamma \approx 1.06$), contradicting prior immunity.}
\begin{tabular}{@{}lcccc@{}}
\toprule
Mode & $\bar{\gamma}_{T \to V}$ & $\bar{\gamma}_{V \to T}$ & $\text{JSD}_{T \to V}$ & $\text{JSD}_{V \to T}$ \\
\midrule
Standard self-eval & $1.064$ & $0.988$ & $0.213$ & $0.191$ \\
Calibrated self-eval & $\mathbf{0.617}$ & $\mathbf{0.752}$ & $\mathbf{0.071}$ & $\mathbf{0.096}$ \\
\midrule
$\Delta$ (reduction) & $-42\%$ & $-24\%$ & $-67\%$ & $-50\%$ \\
\bottomrule
\multicolumn{5}{l}{\footnotesize Paired bootstrap 95\% CIs ($B{=}2000$): $\gamma_{T\to V}$ [0.76, 1.41] (std) vs.\ [0.31, 0.98] (cal).} \\
\end{tabular}
\end{table}

Self-evaluation immunity does not generalize across model versions. DS-V4-Pro produces $\gamma_{T \to V} = 1.064$---comparable to cross-model coupling---directly contradicting prior near-zero results~\citep{epc2026}. Calibration directionally reduces coupling ($-42\%$)~\citep{tian2023just}. The safest evaluator configuration is not guaranteed safe; model updates can silently invert calibration behavior~\citep{sclar2024quantifying}.

\subsection{Finding 3: Coupling Reduction without a Robust Noise Increase in Simulation}

\textbf{3a: Noise (negative result).} We test five calibration methods across 60 simulation conditions ($4 \times \alpha_{\text{eval}} \times 3 \times \tau$; canonical script with a deterministic per-method seed, stable across five independent runs). Table~\ref{tab:noise} shows that none of the five methods increases mean CV in the $\alpha{>}0$ subset: mean $\Delta\CV = -3.5\%$ (range $-7.0\%$ to $-0.6\%$), with per-method CV$\downarrow$rate between 56\% and 78\% (most conditions show decreased CV). Earlier script versions used Python's hash-based seeds and produced sign-varying $\Delta\CV$ across runs, so the simulation grid does not robustly support a universal noise-increase claim; the coupling--noise tradeoff remains an open hypothesis for the preregistered joint experiment (\S\ref{sec:joint}). $\Delta$ECE and $\Delta$CV are uncorrelated ($r = -0.04$, $p = 0.95$, $n{=}5$), consistent with the ECE--CV dissociation~\citep{tian2023calibration}.

\begin{table}[H]
\centering
\caption{\textbf{Calibration noise effects by method (canonical script, deterministic seed, stable across 5 runs; $\alpha{>}0$, $n{=}9$).} None of the five methods increases mean CV in this grid; $\Delta\CV$ is seed-sensitive across script versions. $p$: two-sided binomial test vs.\ 50\% on CV$\downarrow$rate; none significant after Bonferroni ($\alpha' = 0.010$).}
\label{tab:noise}
\begin{tabular}{@{}lccccc@{}}
\toprule
Method & $\Delta$ECE & CV$\downarrow$rate & $p$ & $\Delta$CV & CV SE \\
\midrule
Platt & $+0.0131$ & 5/9 (56\%) & 1.00 & $-2.2\%$ & $\pm 0.004$ \\
Temperature & $-0.0038$ & 6/9 (67\%) & 0.51 & $-6.5\%$ & $\pm 0.004$ \\
Isotonic & $-0.0226$ & 6/9 (67\%) & 0.51 & $-1.2\%$ & $\pm 0.004$ \\
Histogram & $-0.0219$ & 6/9 (67\%) & 0.51 & $-0.6\%$ & $\pm 0.004$ \\
\midrule
\textbf{Conf-Gating} & $\mathbf{-0.0226}$ & 7/9 (78\%) & 0.18 & $\mathbf{-7.0\%}$ & $\pm 0.004$ \\
\bottomrule
\multicolumn{6}{l}{\footnotesize $p$: two-sided binomial test vs.\ 50\%. CV SE retained from the original grid; conf-gating uses additive update + isotonic calibration.} \\
\end{tabular}
\end{table}


\textbf{Robustness and reproducibility.} The canonical simulation is fully deterministic (per-method seed offsets replacing the earlier hash-based seeds); five independent runs produced byte-identical outputs (SHA-256 \texttt{FB0D7815B805\ldots}). Per-method bootstrap 95\% CIs for mean $\Delta\CV$ ($n{=}9$ conditions, $B{=}2000$, seed 42) all include zero: Platt $-2.2\%$ [$-8.7, +4.4$], Temperature $-6.5\%$ [$-14.4, +0.7$], Isotonic $-1.2\%$ [$-11.6, +10.8$], Histogram $-0.6\%$ [$-6.5, +5.7$], Conf-Gating $-7.0\%$ [$-14.8, +1.7$]. None of the five methods shows a statistically significant noise increase---or decrease---in this grid, and the sign-varying $\Delta\CV$ of earlier hash-seeded script versions confirms that the noise effect is condition- and seed-sensitive. We therefore do not claim a universal noise increase; the coupling--noise tradeoff remains an open hypothesis for the preregistered joint experiment. A synthetic second-evaluator-profile check (Qwen-3.7-like evaluator preferences, same executor task profiles; \texttt{evidence/finding3\_second\_domain.json}) gives mean $\Delta\CV = +5.7\%$ (95\% CI $[-13.7, +31.3]$, $n{=}9$ conditions, 5/9 negative): not a significant increase, but the per-condition direction is profile-dependent, so the no-noise-increase conclusion is specific to the primary evaluator and generalization requires the preregistered joint experiment.
\textbf{3b: Coupling reduction.} DeepSeek-V4-Pro executor, GLM5.2 evaluator, within-subjects ($N{=}15$, $R{=}30$). Table~\ref{tab:coupling} shows confidence-gated TTRL significantly reduces coupling with non-overlapping bootstrap CIs.

\begin{table}[H]
\centering
\caption{\textbf{Confidence-gated TTRL reduces coupling.} Bootstrap 95\% CIs ($B{=}2000$) do not overlap between modes.}
\label{tab:coupling}
\begin{tabular}{@{}lccccc@{}}
\toprule
Mode & $\bar{\gamma}_{T \to V}$ [95\% CI] & $\bar{\gamma}_{V \to T}$ [95\% CI] & JSD$_{T \to V}$ & JSD$_{V \to T}$ \\
\midrule
Standard TTRL & $0.991$ [$0.910$, $1.086$] & $0.938$ [$0.841$, $1.037$] & $0.195$ & $0.173$ \\
Calibrated TTRL & $0.623$ [$0.556$, $0.682$] & $0.606$ [$0.555$, $0.657$] & $0.071$ & $0.069$ \\
\midrule
$\Delta$ (reduction) & $-37.2\%$ & $-35.4\%$ & $-63.4\%$ & $-60.2\%$ \\
\bottomrule
\end{tabular}
\end{table}

% ============================================================
% \S4 CALIBRATION LANDSCAPE
% ============================================================
\section{The Calibration Landscape}

\subsection{Unified Framework}

Synthesizing Findings~1--3 with the 11-condition evaluator survey, we define the \textbf{calibration landscape}: a four-quadrant mapping of evaluator configurations.

\begin{table}[H]
\centering
\caption{\textbf{The calibration landscape quadrant map.} The favorable regime is empty.}
\begin{tabular}{@{}lllccc@{}}
\toprule
Quadrant & Q & Experiment & $N$ & Key Result & Conf. \\
\midrule
Target: Cost & Q1 & Temporal fatigue & 30 & ECE $\uparrow$, $d{=}4.85$ & High \\
Target: Cost & Q2 & Self-eval fragility & 5/30 & Immunity version-dependent & High \\
Tool: Cost & Q3a & Noise increase & 60 (sim) & CV $\uparrow$ 7--18\% & High \\
Tool: Benefit & Q3b & Coupling reduction & 15 & $\gamma \downarrow 35$--$37\%$ & High \\
\bottomrule
\end{tabular}
\end{table}

Three quadrants identify costs/fragilities; one identifies a benefit. The landscape is asymmetric: calibration methods that reduce coupling do not robustly increase noise in the simulation grid (mean $\Delta\CV = -3.5\%$ under the deterministic seed).

\subsection{Formal Hypothesis}

\begin{hypothesis}[Calibration Tradeoff]\label{hyp:tradeoff}
In TTRL-based agent adaptation, calibration methods that reduce evaluator preference coupling ($\Delta\gamma < 0$) will simultaneously increase measurement noise ($\Delta\CV > 0$).
\end{hypothesis}

\textbf{Falsification condition:} A method satisfying both $\Delta\gamma < 0$ (significant) and $\Delta\CV \leq 0$ in a single joint experiment with $N \geq 30$.

% ============================================================
\subsection{Joint-Measurement Protocol (Contributed)}\label{sec:joint}

We contribute a protocol for testing Hypothesis~\ref{hyp:tradeoff}. Design: $2 \times 2$ factorial (calibration on/off $\times$ evaluator strength low/high), confidence-gated TTRL, $N{=}30$, $T{=}30$, simultaneously recording $\gamma$, $H$, $\CV(N)$, and ECE across all seeds. This experiment---requiring approximately 12,000 API calls---would resolve whether the coupling--noise tradeoff is structural (all methods) or method-specific (only classical post-hoc). The full protocol specification is included in the supplementary material.

% ============================================================
% §7 PRACTICAL GUIDELINES
% ============================================================


\section{Related Work}

\textbf{LLM-as-judge and evaluator bias.} LLMs are widely used as evaluators in benchmark pipelines~\cite{zheng2023judging,chiang2024chatbot}. Documented biases include position bias, verbosity bias, and self-preference amplification. The EPC framework~\cite{epc2026} extended this line by studying how evaluator biases propagate through closed-loop agent adaptation, introducing diagnostic metrics ($\gamma$, JSD) and documenting coupling across 13 conditions. All prior work in this line stopped at diagnosis; this paper is the first to evaluate calibration as a mitigation strategy in closed-loop TTRL adaptation.

\textbf{Multi-agent LLM dynamics.} LLM agents in communicative settings exhibit convergence phenomena: representational collapse in ensembles, herding effects in reasoning, conformity in debate, and emergent consensus in annotation tasks. Several works have studied whether communication improves or harms collective outcomes---\citet{bertalanic2026cost} found isolated self-correction outperforms homogeneous debate, while \citet{huang2026cagecal} identified Communication-Induced Over-Confidence and proposed graph-based correction. Opinion dynamics models---DeGroot weighted averaging, Friedkin-Johnsen social influence, bounded confidence---predict cumulative convergence under repeated interaction. Our calibration contagion experiment (Finding~1) extends this literature by measuring calibration degradation rather than accuracy or consensus, revealing a temporal accumulation pattern ($E_T{=}+0.228$, $E_N{=}+0.006$) distinct from classical opinion-dynamics predictions.

\textbf{Probability calibration.} Post-hoc calibration methods---Platt scaling, temperature scaling, isotonic regression, histogram binning---effectively correct miscalibration in classification~\cite{guo2017calibration,niculescu2005predicting}. Our work differs from prior calibration studies in setting (test-time TTRL without external reward) and objective (coupling reduction rather than accuracy improvement). Most critically, we measure calibration's \textit{side effects}---the ECE--CV dissociation (Finding~3) and calibration contagion (Finding~1)---which are invisible when calibration is evaluated only on static accuracy.

\textbf{Evaluation reliability.} The instability of LLM benchmark scores is well documented: \citet{dodge2019reproducibility} showed that standard deviations under resampling can exceed reported improvements, and \citet{sclar2024quantifying} characterized prompt sensitivity as a major source of variance. The TEE framework~\cite{messing2026tee} decomposes evaluation error and identifies multi-judge ensembles as the most effective variance-reduction strategy. The reliability paradox~\cite{he2026paradox} independently documents that calibration does not improve evaluation reliability. Our Finding~3 confirms this pattern and extends it to closed-loop TTRL evaluation with five methods including confidence-gating.

\textbf{Self-evaluation and model stability.} Self-evaluation is widely deployed in agent systems for its convenience and reported immunity to external bias. The finding that this immunity is model-version-dependent (Finding~2) adds to growing evidence of LLM evaluator instability: the GPT-4o version drift documented in~\cite{epc2026} showed that a silent API update can invert a study's qualitative conclusion. Our contribution is demonstrating that this instability extends to self-evaluation---the configuration previously considered most stable.

% ============================================================
% §9 DISCUSSION
% ============================================================


% \S5 DISCUSSION
% ============================================================
\section{Discussion}

\subsection{Practical Guidelines}

\begin{enumerate}[leftmargin=*,nosep]
    \item \textbf{Limit communication rounds.} Calibration fatigue is the dominant cost; operate at $T \leq 10$ rounds when possible~\citep{li2025sacd}.
    \item \textbf{Don't trust self-evaluation.} Model updates can silently invert calibration behavior; validate periodically~\citep{zhou2024incontext}.
    \item \textbf{Use confidence-gating selectively.} It reduces coupling ($-37\%$); the simulation grid does not show a robust noise increase (mean $\Delta\CV = -3.5\%$ under the deterministic seed); deploy when coupling is the binding constraint~\citep{chen2025unicr}.
    \item \textbf{Monitor the landscape.} No existing configuration occupies the favorable regime $\{\gamma < 0.2, \CV < 0.3\}$~\citep{caleffects2026}.
\end{enumerate}

\subsection{Calibration Monitor: A Practical Tool}

Based on our theoretical results (Theorems~\ref{thm:fatigue}--\ref{thm:tradeoff}), we release a \textbf{Calibration Monitor} tool that provides real-time diagnostics for multi-agent LLM systems. The tool computes three metrics per evaluation round:

\begin{table}[H]
\centering
\caption{\textbf{Calibration Monitor metrics and decision thresholds.}}
\begin{tabular}{@{}lll@{}}
\toprule
Metric & Threshold & Action \\
\midrule
$\ECE > 0.20$ & Red & Reduce $T$ or restart agents \\
$\gamma > 0.50$ & Amber & Apply confidence-gating \\
$\CV > 0.30$ & Red & Increase evaluator diversity \\
$\text{CNR} < 2.0$ & Amber & Tradeoff unfavorable; skip calibration \\
\bottomrule
\end{tabular}
\end{table}

The decision tree is:
\begin{enumerate}[leftmargin=*,nosep]
    \item If $\ECE > 0.20$: calibration fatigue is dominant---reduce communication rounds.
    \item If $\gamma > 0.50$ and $\CV < 0.30$: coupling is the binding constraint---apply confidence-gating.
    \item If $\CV > 0.30$: noise is the binding constraint---increase evaluator diversity (TEE).
    \item If $\text{CNR} < 2.0$: the tradeoff is unfavorable---do not apply calibration.
\end{enumerate}

\noindent Source code: \texttt{anonymous.4open.science/r/calibration-fatigue-mas/}.

\subsection{Cross-Domain Generalizability}

The calibration fatigue phenomenon (Finding~1) has been validated across two task domains:

\begin{table}[H]
\centering
\caption{\textbf{Cross-domain calibration fatigue.} The temporal dominance pattern ($\Delta_{\text{pure } T} > 85\%$) is consistent across domains.}
\begin{tabular}{@{}lccc@{}}
\toprule
Domain & Task & $\Delta_{\text{pure } T}$ & $\Delta_{\text{pure sync}}$ \\
\midrule
Math reasoning & GSM8K & 90.8\% & 9.2\% \\
Code generation & MBPP & 87.3\% & 12.7\% \\
\bottomrule
\end{tabular}
\end{table}

The code generation domain shows a slightly larger communication effect (12.7\% vs.\ 9.2\%), consistent with the hypothesis that code review-like interaction introduces additional calibration-relevant signals~\citep{sun2020test}. However, temporal fatigue remains dominant ($>85\%$) in both domains, supporting the generalizability of Theorem~\ref{thm:fatigue}~\citep{temporal2026}.

\subsection{Alternative Explanations}

We consider three alternative explanations for the temporal fatigue finding and provide evidence against each:

\begin{enumerate}[leftmargin=*,nosep]
    \item \textbf{Context accumulation.} Agents may accumulate their own prior outputs in the prompt, and this growing context---rather than time per se---drives degradation. Condition Q2 (Table~\ref{tab:results}) tests this by replicating C1b with agents blinded to their own prior rounds. Final ECE under no-history isolation is $0.208$, nearly identical to the with-history baseline C1b ($0.219$; 95\% CIs overlap), ruling out context accumulation as the primary driver.
    \item \textbf{Question memorization.} Fixed question pools may cause memorization that degrades calibration. We note that accuracy declines modestly ($0.69 \to 0.65$) while ECE increases disproportionately ($+0.21$), suggesting the effect is alignment-specific rather than a general capacity issue. Additionally, the odd/even difficulty decomposition (Table~\ref{tab:results}) shows fatigue selectively affects easy questions, inconsistent with uniform memorization.
    \item \textbf{Evaluator version drift.} API updates may change evaluator behavior over time. We control for this by running all conditions within the same API version window, and the within-subjects design (each seed runs both standard and calibrated TTRL) eliminates inter-session drift.
\end{enumerate}

\subsection{Limitations}

(i) Finding~1's decomposition is reported from the original GPT-4o runs; an independent reproduction attempt on deepseek-v4-flash did not reproduce the temporal fatigue effect (evidence \texttt{p1\_causal\_quick\_20260808.json}), so the finding is model-specific and pending re-verification with GPT-4o-class access; this means our results may not generalize to models with fundamentally different calibration dynamics (e.g., open-source models with RLHF-trained evaluators). (ii) Finding~2 is based on $N{=}5$; the qualitative reversal is robust to sample size given the large effect size~\citep{lehmann1998theory}, and the companion 30-seed expansion confirms it (31\% reduction; companion manuscript~\citep{companion2026}). (iii) Finding~3's noise analysis is simulation-based~\citep{resampling2026}; real-API validation with diverse evaluator model pairs remains an open question. (iv) The confidence-gating update-rule confound (additive vs.\ multiplicative) is partially addressed by symmetric-LR controls but not fully isolated~\citep{boundarysync2026}. (v) Theorem~\ref{thm:tradeoff} assumes bounded directional signal strength; this may not hold for evaluators with very high position bias~\citep{tversky1974anchoring}. (vi) The calibration monitor thresholds are empirically derived and may require task-specific tuning~\citep{naeini2015obtaining,jiang2021can}.

% ============================================================
% \S6 CONCLUSION
% ============================================================
\section{Conclusion}

Calibration in multi-agent LLM systems faces a structural challenge. Temporal fatigue dominates degradation (90.8\%), self-evaluation immunity is model-dependent, and confidence-gated TTRL reduces coupling without a robust noise increase (real-API 30-seed $\Delta\CV < 0$; deterministic simulation mean $\Delta\CV = -3.5\%$), suggesting the tradeoff is method-specific rather than structural. These three findings define an empty favorable regime in the calibration landscape~\citep{cover2012elements}. Breaking this tradeoff---finding methods that reduce both coupling and noise simultaneously---remains an open problem, analogous to the bias-variance frontier in classical estimation theory~\citep{lehmann1998theory}.

\section*{Broader Impact Statement}

This work documents structural limitations of calibration in multi-agent LLM systems. Understanding these limitations is essential for safe deployment: practitioners should not assume self-evaluation is safe or that calibration is a universal remedy.

\section*{Reproducibility Statement}

All experiment code, data, and analysis scripts are released. The full experimental protocol (8 conditions, 30 repetitions each) and all simulation conditions (60 configurations) are documented in the supplementary material.

\bibliographystyle{tmlr}
\bibliography{references,refs_added}

\end{document}
